\section{Attestation、boot、lifecycle}

\begin{frame}[t,shrink=6]{Attestation、boot、lifecycle：方向开场}
\DeckKicker{方向开场}
\SlideMeta{证据等级}{方向综述}{来源状态：公开材料综合}
{\large\bfseries\color{Ink} Attestation、boot、lifecycle 的核心是先界定保护对象、管理边界和证据强度，再用三篇主讲材料串起技术演进。\par}\vspace{1.2mm}
\begin{columns}[T,totalwidth=\textwidth]
  \begin{column}{0.45\textwidth}
    \fcolorbox{Rule}{Paper}{%
  \begin{minipage}[t]{0.97\linewidth}
    \vspace{1.1mm}
    \hspace{1.4mm}\begin{minipage}{0.92\linewidth}
      \PanelTitle{内容说明}
      \scriptsize \begin{itemize}
  \item 区分 secure boot、launch attestation、runtime/control-flow evidence、evidence profile 与 verifier policy。
  \item 固定三篇主讲材料；`rats\_rfc` 主讲 attestation architecture，`eat\_rfc` 作为 token-format auxiliary；C-FLAT/LO-FAT/PDRIMA 作为 runtime-attestation lineage auxiliary。
  \item 先讲方向痛点和设计轴，再逐篇说明贡献、限制、证据强度和可落地条件。
  \item 对规范、Survey、draft、vendor evidence 保持显式标注，避免把背景材料写成一手系统证明。
\end{itemize}
    \end{minipage}
    \vspace{1mm}
  \end{minipage}}
  \end{column}
  \begin{column}{0.49\textwidth}
    \fcolorbox{Rule}{Panel}{%
  \begin{minipage}[t]{0.97\linewidth}
    \vspace{1.1mm}
    \hspace{1.4mm}\begin{minipage}{0.92\linewidth}
      \PanelTitle{三篇主讲材料的分工}
      \scriptsize {\tiny
\begin{tabularx}{\linewidth}{@{}YYY@{}}
\toprule
\textbf{论文} & \textbf{定位} & \textbf{证据边界} \\
\midrule
SWATT: 软件 attestation 早期代表 & 开创/基础入口 & E1 / 基础证据 \\
RATS 架构 & 代表性改进 & E0 / Spec-standard SOTA \\
VRASED verified RA co-design & 当前 SOTA/补充边界 & E1 / 同行评审改进证据 \\
\bottomrule
\end{tabularx}
}
    \end{minipage}
    \vspace{1mm}
  \end{minipage}}
  \end{column}
\end{columns}
\SourceFootnote{来源：论文原文、官方规范或公开材料｜证据等级：见本页 badge}
\end{frame}


\begin{frame}[t,shrink=6]{SWATT: 软件 attestation 早期代表：内容摘要}
\DeckKicker{内容摘要}
\SlideMeta{证据等级}{E1 / 基础证据}{来源状态：本地 PDF 已核验}
{\large\bfseries\color{Ink} 提出软件 attestation 的早期代表方法，用 timing-sensitive checksum 让远程 verifier 检查设备内存内容\par}\vspace{1.2mm}
\begin{columns}[T,totalwidth=\textwidth]
  \begin{column}{0.45\textwidth}
    \fcolorbox{Rule}{Paper}{%
  \begin{minipage}[t]{0.97\linewidth}
    \vspace{1.1mm}
    \hspace{1.4mm}\begin{minipage}{0.92\linewidth}
      \PanelTitle{内容说明}
      \scriptsize \begin{itemize}
  \item 提出软件 attestation 的早期代表方法，用 timing-sensitive checksum 让远程 verifier 检查设备内存内容
  \item SWATT is the classic starting point for remote software attestation problem framing.
\end{itemize}
    \end{minipage}
    \vspace{1mm}
  \end{minipage}}
  \end{column}
  \begin{column}{0.49\textwidth}
    \fcolorbox{Rule}{Panel}{%
  \begin{minipage}[t]{0.97\linewidth}
    \vspace{1.1mm}
    \hspace{1.4mm}\begin{minipage}{0.92\linewidth}
      \PanelTitle{证据定位}
      \scriptsize {\tiny
\begin{tabularx}{\linewidth}{@{}YY@{}}
\toprule
\textbf{字段} & \textbf{内容} \\
\midrule
论文定位 & 开创/基础入口 \\
材料类型 & 系统论文 \\
证据等级 & E1 / 基础证据 \\
来源状态 & 本地 PDF 已核验 \\
\bottomrule
\end{tabularx}
}
    \end{minipage}
    \vspace{1mm}
  \end{minipage}}
  \end{column}
\end{columns}
\SourceFootnote{来源：SWATT: 软件 attestation 早期代表｜证据等级：E1 / 基础证据}
\end{frame}


\begin{frame}[t,shrink=6]{SWATT: 软件 attestation 早期代表：研究背景}
\DeckKicker{研究背景}
\SlideMeta{证据等级}{E1 / 基础证据}{来源状态：本地 PDF 已核验}
{\large\bfseries\color{Ink} 低端嵌入式设备缺少硬件 root of trust，但仍需要远程确认代码完整性和启动状态\par}\vspace{1.2mm}
\begin{columns}[T,totalwidth=\textwidth]
  \begin{column}{0.45\textwidth}
    \fcolorbox{Rule}{Paper}{%
  \begin{minipage}[t]{0.97\linewidth}
    \vspace{1.1mm}
    \hspace{1.4mm}\begin{minipage}{0.92\linewidth}
      \PanelTitle{内容说明}
      \scriptsize \begin{itemize}
  \item 低端嵌入式设备缺少硬件 root of trust，但仍需要远程确认代码完整性和启动状态
\end{itemize}
    \end{minipage}
    \vspace{1mm}
  \end{minipage}}
  \end{column}
  \begin{column}{0.49\textwidth}
    \fcolorbox{Rule}{Panel}{%
  \begin{minipage}[t]{0.97\linewidth}
    \vspace{1.1mm}
    \hspace{1.4mm}\begin{minipage}{0.92\linewidth}
      \PanelTitle{问题边界}
      \scriptsize {\tiny
\begin{tabularx}{\linewidth}{@{}YY@{}}
\toprule
\textbf{维度} & \textbf{说明} \\
\midrule
保护对象 & Attestation、boot、lifecycle \\
传统瓶颈 & 区分 secure boot、launch attestation、runtime/control-flow evidence... \\
本文入口 & SWATT: 软件 attestation 早期代表 \\
证据限制 & E1 / 基础证据 \\
\bottomrule
\end{tabularx}
}
    \end{minipage}
    \vspace{1mm}
  \end{minipage}}
  \end{column}
\end{columns}
\SourceFootnote{来源：SWATT: 软件 attestation 早期代表｜证据等级：E1 / 基础证据}
\end{frame}


\begin{frame}[t,shrink=6]{SWATT: 软件 attestation 早期代表：解决方案}
\DeckKicker{解决方案}
\SlideMeta{证据等级}{E1 / 基础证据}{来源状态：本地 PDF 已核验}
{\large\bfseries\color{Ink} 用严格时间假设约束 prover 计算能力，使篡改设备难以伪造完整内存 checksum\par}\vspace{1.2mm}
\begin{columns}[T,totalwidth=\textwidth]
  \begin{column}{0.45\textwidth}
    \fcolorbox{Rule}{Paper}{%
  \begin{minipage}[t]{0.97\linewidth}
    \vspace{1.1mm}
    \hspace{1.4mm}\begin{minipage}{0.92\linewidth}
      \PanelTitle{内容说明}
      \scriptsize \begin{itemize}
  \item 用严格时间假设约束 prover 计算能力，使篡改设备难以伪造完整内存 checksum
\end{itemize}
    \end{minipage}
    \vspace{1mm}
  \end{minipage}}
  \end{column}
  \begin{column}{0.49\textwidth}
    \fcolorbox{Rule}{Panel}{%
  \begin{minipage}[t]{0.97\linewidth}
    \vspace{1.1mm}
    \hspace{1.4mm}\begin{minipage}{0.92\linewidth}
      \PanelTitle{核心设计拆解}
      \scriptsize \begin{itemize}
  \item 01. Software-only remote attestation
  \item 02. Timing-sensitive memory checksum
  \item 03. Verifier/prover challenge-response model
\end{itemize}
    \end{minipage}
    \vspace{1mm}
  \end{minipage}}
  \end{column}
\end{columns}
\SourceFootnote{来源：SWATT: 软件 attestation 早期代表｜证据等级：E1 / 基础证据}
\end{frame}


\begin{frame}[t,shrink=6]{SWATT: 软件 attestation 早期代表：实验结果}
\DeckKicker{实验结果}
\SlideMeta{证据等级}{E1 / 基础证据}{来源状态：本地 PDF 已核验}
{\large\bfseries\color{Ink} 论文在嵌入式节点场景评估\par}\vspace{1.2mm}
\begin{columns}[T,totalwidth=\textwidth]
  \begin{column}{0.45\textwidth}
    \fcolorbox{Rule}{Paper}{%
  \begin{minipage}[t]{0.97\linewidth}
    \vspace{1.1mm}
    \hspace{1.4mm}\begin{minipage}{0.92\linewidth}
      \PanelTitle{内容说明}
      \scriptsize \begin{itemize}
  \item 论文在嵌入式节点场景评估
  \item 安全性强依赖 timing、硬件和网络假设
\end{itemize}
    \end{minipage}
    \vspace{1mm}
  \end{minipage}}
  \end{column}
  \begin{column}{0.49\textwidth}
    \fcolorbox{Rule}{Panel}{%
  \begin{minipage}[t]{0.97\linewidth}
    \vspace{1.1mm}
    \hspace{1.4mm}\begin{minipage}{0.92\linewidth}
      \PanelTitle{实验/证据边界}
      \scriptsize {\tiny
\begin{tabularx}{\linewidth}{@{}YYY@{}}
\toprule
\textbf{对象} & \textbf{可支撑} & \textbf{边界} \\
\midrule
数据/来源 & SWATT: 软件 attestation 早期代表 & 本地 PDF 已核验 \\
Baseline/对照 & 以论文或规范原文为准 & 不跨材料外推 \\
指标/对象 & 只使用当前来源可证明的信息 & E1 / 基础证据 \\
使用边界 & 可进入本方向演进图 & 不能替代更强证据 \\
\bottomrule
\end{tabularx}
}
    \end{minipage}
    \vspace{1mm}
  \end{minipage}}
  \end{column}
\end{columns}
\SourceFootnote{来源：SWATT: 软件 attestation 早期代表｜证据等级：E1 / 基础证据}
\end{frame}


\begin{frame}[t,shrink=6]{SWATT: 软件 attestation 早期代表：文章评价}
\DeckKicker{文章评价}
\SlideMeta{证据等级}{E1 / 基础证据}{来源状态：本地 PDF 已核验}
{\large\bfseries\color{Ink} 历史价值高，能解释 attestation 从软件测量走向硬件 RoT 的动因\par}\vspace{1.2mm}
\begin{columns}[T,totalwidth=\textwidth]
  \begin{column}{0.45\textwidth}
    \fcolorbox{Rule}{Paper}{%
  \begin{minipage}[t]{0.97\linewidth}
    \vspace{1.1mm}
    \hspace{1.4mm}\begin{minipage}{0.92\linewidth}
      \PanelTitle{内容说明}
      \scriptsize \begin{itemize}
  \item 设计优势：历史价值高，能解释 attestation 从软件测量走向硬件 RoT 的动因。
  \item 局限性：面对现代强对手和复杂 SoC/TEE 时适用性有限。
  \item 商业化潜力：现代商业方案应转向硬件 RoT、标准化 evidence 和 verifier policy；SWATT 主要作为历史基线。
  \item 证据边界：E1 / Foundational；本地 PDF 已核验。
\end{itemize}
    \end{minipage}
    \vspace{1mm}
  \end{minipage}}
  \end{column}
  \begin{column}{0.49\textwidth}
    \fcolorbox{Rule}{Panel}{%
  \begin{minipage}[t]{0.97\linewidth}
    \vspace{1.1mm}
    \hspace{1.4mm}\begin{minipage}{0.92\linewidth}
      \PanelTitle{文章评价}
      \scriptsize {\tiny
\begin{tabularx}{\linewidth}{@{}YY@{}}
\toprule
\textbf{维度} & \textbf{判断} \\
\midrule
设计优势 & 历史价值高，能解释 attestation 从软件测量走向硬件 RoT 的动因。 \\
主要局限 & 面对现代强对手和复杂 SoC/TEE 时适用性有限。 \\
商业落地 & 现代商业方案应转向硬件 RoT、标准化 evidence 和 verifier policy；SWATT 主要作为历史基线。 \\
讲稿定位 & 开创/基础入口; 证据强度 E1 \\
\bottomrule
\end{tabularx}
}
    \end{minipage}
    \vspace{1mm}
  \end{minipage}}
  \end{column}
\end{columns}
\SourceFootnote{来源：SWATT: 软件 attestation 早期代表｜证据等级：E1 / 基础证据}
\end{frame}


\begin{frame}[t,shrink=6]{RATS 架构：内容摘要}
\DeckKicker{内容摘要}
\SlideMeta{证据等级}{E0 / Spec-standard SOTA}{来源状态：本地 PDF 已核验}
{\large\bfseries\color{Ink} 定义 attester、verifier、relying party、evidence、endorsements、reference values、attestation results 和 appraisal policy\par}\vspace{1.2mm}
\begin{columns}[T,totalwidth=\textwidth]
  \begin{column}{0.45\textwidth}
    \fcolorbox{Rule}{Paper}{%
  \begin{minipage}[t]{0.97\linewidth}
    \vspace{1.1mm}
    \hspace{1.4mm}\begin{minipage}{0.92\linewidth}
      \PanelTitle{内容说明}
      \scriptsize \begin{itemize}
  \item 定义 attester、verifier、relying party、evidence、endorsements、reference values、attestation...
  \item RFC 9334 defines the current RATS architecture for cross-TEE remote attestation.
\end{itemize}
    \end{minipage}
    \vspace{1mm}
  \end{minipage}}
  \end{column}
  \begin{column}{0.49\textwidth}
    \fcolorbox{Rule}{Panel}{%
  \begin{minipage}[t]{0.97\linewidth}
    \vspace{1.1mm}
    \hspace{1.4mm}\begin{minipage}{0.92\linewidth}
      \PanelTitle{证据定位}
      \scriptsize {\tiny
\begin{tabularx}{\linewidth}{@{}YY@{}}
\toprule
\textbf{字段} & \textbf{内容} \\
\midrule
论文定位 & 代表性改进 \\
材料类型 & 规范/标准材料 \\
证据等级 & E0 / Spec-standard SOTA \\
来源状态 & 本地 PDF 已核验 \\
\bottomrule
\end{tabularx}
}
    \end{minipage}
    \vspace{1mm}
  \end{minipage}}
  \end{column}
\end{columns}
\SourceFootnote{来源：RATS 架构｜证据等级：E0 / Spec-standard SOTA}
\end{frame}


\begin{frame}[t,shrink=6]{RATS 架构：研究背景}
\DeckKicker{研究背景}
\SlideMeta{证据等级}{E0 / Spec-standard SOTA}{来源状态：本地 PDF 已核验}
{\large\bfseries\color{Ink} 多平台 attestation 需要统一角色、消息、freshness 和 verifier policy 词汇，否则 CCA、CoVE、DPU/NIC evidence 难以对齐\par}\vspace{1.2mm}
\begin{columns}[T,totalwidth=\textwidth]
  \begin{column}{0.45\textwidth}
    \fcolorbox{Rule}{Paper}{%
  \begin{minipage}[t]{0.97\linewidth}
    \vspace{1.1mm}
    \hspace{1.4mm}\begin{minipage}{0.92\linewidth}
      \PanelTitle{内容说明}
      \scriptsize \begin{itemize}
  \item 多平台 attestation 需要统一角色、消息、freshness 和 verifier policy 词汇，否则 CCA、CoVE、DPU/NIC evidence 难以对齐
\end{itemize}
    \end{minipage}
    \vspace{1mm}
  \end{minipage}}
  \end{column}
  \begin{column}{0.49\textwidth}
    \fcolorbox{Rule}{Panel}{%
  \begin{minipage}[t]{0.97\linewidth}
    \vspace{1.1mm}
    \hspace{1.4mm}\begin{minipage}{0.92\linewidth}
      \PanelTitle{问题边界}
      \scriptsize {\tiny
\begin{tabularx}{\linewidth}{@{}YY@{}}
\toprule
\textbf{维度} & \textbf{说明} \\
\midrule
保护对象 & Attestation、boot、lifecycle \\
传统瓶颈 & 区分 secure boot、launch attestation、runtime/control-flow evidence... \\
本文入口 & RATS 架构 \\
证据限制 & E0 / Spec-standard SOTA \\
\bottomrule
\end{tabularx}
}
    \end{minipage}
    \vspace{1mm}
  \end{minipage}}
  \end{column}
\end{columns}
\SourceFootnote{来源：RATS 架构｜证据等级：E0 / Spec-standard SOTA}
\end{frame}


\begin{frame}[t,shrink=6]{RATS 架构：解决方案}
\DeckKicker{解决方案}
\SlideMeta{证据等级}{E0 / Spec-standard SOTA}{来源状态：本地 PDF 已核验}
{\large\bfseries\color{Ink} RATS 规范化 evidence 生成、传递、评估角色与拓扑\par}\vspace{1.2mm}
\begin{columns}[T,totalwidth=\textwidth]
  \begin{column}{0.45\textwidth}
    \fcolorbox{Rule}{Paper}{%
  \begin{minipage}[t]{0.97\linewidth}
    \vspace{1.1mm}
    \hspace{1.4mm}\begin{minipage}{0.92\linewidth}
      \PanelTitle{内容说明}
      \scriptsize \begin{itemize}
  \item RATS 规范化 evidence 生成、传递、评估角色与拓扑
  \item `eat\_rfc` 辅助说明实体状态和安全属性 claims 的 token 编码
\end{itemize}
    \end{minipage}
    \vspace{1mm}
  \end{minipage}}
  \end{column}
  \begin{column}{0.49\textwidth}
    \fcolorbox{Rule}{Panel}{%
  \begin{minipage}[t]{0.97\linewidth}
    \vspace{1.1mm}
    \hspace{1.4mm}\begin{minipage}{0.92\linewidth}
      \PanelTitle{核心设计拆解}
      \scriptsize \begin{itemize}
  \item 01. Attester/verifier/relying-party roles
  \item 02. Evidence, endorsements, reference values and appraisal policy
  \item 03. Attestation results and freshness vocabulary
\end{itemize}
    \end{minipage}
    \vspace{1mm}
  \end{minipage}}
  \end{column}
\end{columns}
\SourceFootnote{来源：RATS 架构｜证据等级：E0 / Spec-standard SOTA}
\end{frame}


\begin{frame}[t,shrink=6]{RATS 架构：实验结果}
\DeckKicker{实验结果}
\SlideMeta{证据等级}{E0 / Spec-standard SOTA}{来源状态：本地 PDF 已核验}
{\large\bfseries\color{Ink} 规范/Survey，无新实验；RATS 架构 只支撑术语、taxonomy、接口语义或证据边界。\par}\vspace{1.2mm}
\begin{columns}[T,totalwidth=\textwidth]
  \begin{column}{0.45\textwidth}
    \fcolorbox{Rule}{Paper}{%
  \begin{minipage}[t]{0.97\linewidth}
    \vspace{1.1mm}
    \hspace{1.4mm}\begin{minipage}{0.92\linewidth}
      \PanelTitle{内容说明}
      \scriptsize \begin{itemize}
  \item 本页不展示独立系统实验，也不把二手归纳写成一手结果。
  \item 可支撑：RFC 9334 defines the current RATS architecture for cross-TEE remote attestation.
  \item 证据等级：E0 / Spec-standard SOTA；来源状态：本地 PDF 已核验。
  \item 涉及具体平台、协议状态或数值时，转到原始论文、官方规范或厂商材料核验。
\end{itemize}
    \end{minipage}
    \vspace{1mm}
  \end{minipage}}
  \end{column}
  \begin{column}{0.49\textwidth}
    \fcolorbox{Rule}{Panel}{%
  \begin{minipage}[t]{0.97\linewidth}
    \vspace{1.1mm}
    \hspace{1.4mm}\begin{minipage}{0.92\linewidth}
      \PanelTitle{实验/证据边界}
      \scriptsize {\tiny
\begin{tabularx}{\linewidth}{@{}YYY@{}}
\toprule
\textbf{对象} & \textbf{可支撑} & \textbf{边界} \\
\midrule
数据/来源 & RATS 架构 & 本地 PDF 已核验 \\
Baseline/对照 & 以论文或规范原文为准 & 不跨材料外推 \\
指标/对象 & 只使用当前来源可证明的信息 & E0 / Spec-standard SOTA \\
使用边界 & 可进入本方向演进图 & 不能替代更强证据 \\
\bottomrule
\end{tabularx}
}
    \end{minipage}
    \vspace{1mm}
  \end{minipage}}
  \end{column}
\end{columns}
\SourceFootnote{来源：RATS 架构｜证据等级：E0 / Spec-standard SOTA}
\end{frame}


\begin{frame}[t,shrink=6]{RATS 架构：文章评价}
\DeckKicker{文章评价}
\SlideMeta{证据等级}{E0 / Spec-standard SOTA}{来源状态：本地 PDF 已核验}
{\large\bfseries\color{Ink} 评价：RATS 架构 适合做 taxonomy/规范入口；规范/Survey，无新实验，不能替代一手系统证据。\par}\vspace{1.2mm}
\begin{columns}[T,totalwidth=\textwidth]
  \begin{column}{0.45\textwidth}
    \fcolorbox{Rule}{Paper}{%
  \begin{minipage}[t]{0.97\linewidth}
    \vspace{1.1mm}
    \hspace{1.4mm}\begin{minipage}{0.92\linewidth}
      \PanelTitle{内容说明}
      \scriptsize \begin{itemize}
  \item 设计优势：是跨平台 attestation 术语和 verifier policy 的权威基线。
  \item 局限性：不证明任一平台测量链充分，也不定义完整 token claim encoding。
  \item 商业化潜力：标准化落地潜力强；风险在平台 evidence profile、endorsement 分发和 verifier policy 运维复杂。
  \item 规范/Survey，无新实验；具体平台结论转到一手材料核验。
\end{itemize}
    \end{minipage}
    \vspace{1mm}
  \end{minipage}}
  \end{column}
  \begin{column}{0.49\textwidth}
    \fcolorbox{Rule}{Panel}{%
  \begin{minipage}[t]{0.97\linewidth}
    \vspace{1.1mm}
    \hspace{1.4mm}\begin{minipage}{0.92\linewidth}
      \PanelTitle{文章评价}
      \scriptsize {\tiny
\begin{tabularx}{\linewidth}{@{}YY@{}}
\toprule
\textbf{维度} & \textbf{判断} \\
\midrule
设计优势 & 是跨平台 attestation 术语和 verifier policy 的权威基线。 \\
主要局限 & 不证明任一平台测量链充分，也不定义完整 token claim encoding。 \\
商业落地 & 标准化落地潜力强；风险在平台 evidence profile、endorsement 分发和 verifier policy 运维复杂。 \\
讲稿定位 & 代表性改进; 证据强度 E0 \\
\bottomrule
\end{tabularx}
}
    \end{minipage}
    \vspace{1mm}
  \end{minipage}}
  \end{column}
\end{columns}
\SourceFootnote{来源：RATS 架构｜证据等级：E0 / Spec-standard SOTA}
\end{frame}


\begin{frame}[t,shrink=6]{VRASED verified RA co-design：内容摘要}
\DeckKicker{内容摘要}
\SlideMeta{证据等级}{E1 / 同行评审改进证据}{来源状态：本地 PDF 已核验}
{\large\bfseries\color{Ink} 设计并验证面向 simple embedded devices 的 hardware/software remote attestation root\par}\vspace{1.2mm}
\begin{columns}[T,totalwidth=\textwidth]
  \begin{column}{0.45\textwidth}
    \fcolorbox{Rule}{Paper}{%
  \begin{minipage}[t]{0.97\linewidth}
    \vspace{1.1mm}
    \hspace{1.4mm}\begin{minipage}{0.92\linewidth}
      \PanelTitle{内容说明}
      \scriptsize \begin{itemize}
  \item 设计并验证面向 simple embedded devices 的 hardware/software remote attestation root
  \item VRASED is the strongest peer-reviewed verified RA co-design anchor for hardware/software...
\end{itemize}
    \end{minipage}
    \vspace{1mm}
  \end{minipage}}
  \end{column}
  \begin{column}{0.49\textwidth}
    \fcolorbox{Rule}{Panel}{%
  \begin{minipage}[t]{0.97\linewidth}
    \vspace{1.1mm}
    \hspace{1.4mm}\begin{minipage}{0.92\linewidth}
      \PanelTitle{证据定位}
      \scriptsize {\tiny
\begin{tabularx}{\linewidth}{@{}YY@{}}
\toprule
\textbf{字段} & \textbf{内容} \\
\midrule
论文定位 & 当前 SOTA/补充边界 \\
材料类型 & 系统论文 \\
证据等级 & E1 / 同行评审改进证据 \\
来源状态 & 本地 PDF 已核验 \\
\bottomrule
\end{tabularx}
}
    \end{minipage}
    \vspace{1mm}
  \end{minipage}}
  \end{column}
\end{columns}
\SourceFootnote{来源：VRASED verified RA co-design｜证据等级：E1 / 同行评审改进证据}
\end{frame}


\begin{frame}[t,shrink=6]{VRASED verified RA co-design：研究背景}
\DeckKicker{研究背景}
\SlideMeta{证据等级}{E1 / 同行评审改进证据}{来源状态：本地 PDF 已核验}
{\large\bfseries\color{Ink} RA 安全性依赖硬件隔离、软件 routine 和协议共同成立，但早期方案缺少实现级机器可检查证明\par}\vspace{1.2mm}
\begin{columns}[T,totalwidth=\textwidth]
  \begin{column}{0.45\textwidth}
    \fcolorbox{Rule}{Paper}{%
  \begin{minipage}[t]{0.97\linewidth}
    \vspace{1.1mm}
    \hspace{1.4mm}\begin{minipage}{0.92\linewidth}
      \PanelTitle{内容说明}
      \scriptsize \begin{itemize}
  \item RA 安全性依赖硬件隔离、软件 routine 和协议共同成立，但早期方案缺少实现级机器可检查证明
\end{itemize}
    \end{minipage}
    \vspace{1mm}
  \end{minipage}}
  \end{column}
  \begin{column}{0.49\textwidth}
    \fcolorbox{Rule}{Panel}{%
  \begin{minipage}[t]{0.97\linewidth}
    \vspace{1.1mm}
    \hspace{1.4mm}\begin{minipage}{0.92\linewidth}
      \PanelTitle{问题边界}
      \scriptsize {\tiny
\begin{tabularx}{\linewidth}{@{}YY@{}}
\toprule
\textbf{维度} & \textbf{说明} \\
\midrule
保护对象 & Attestation、boot、lifecycle \\
传统瓶颈 & 区分 secure boot、launch attestation、runtime/control-flow evidence... \\
本文入口 & VRASED verified RA co-design \\
证据限制 & E1 / 同行评审改进证据 \\
\bottomrule
\end{tabularx}
}
    \end{minipage}
    \vspace{1mm}
  \end{minipage}}
  \end{column}
\end{columns}
\SourceFootnote{来源：VRASED verified RA co-design｜证据等级：E1 / 同行评审改进证据}
\end{frame}


\begin{frame}[t,shrink=6]{VRASED verified RA co-design：解决方案}
\DeckKicker{解决方案}
\SlideMeta{证据等级}{E1 / 同行评审改进证据}{来源状态：本地 PDF 已核验}
{\large\bfseries\color{Ink} 用少量硬件扩展保护 attestation 代码、密钥和内存访问，并形式化验证整体 RA 安全属性\par}\vspace{1.2mm}
\begin{columns}[T,totalwidth=\textwidth]
  \begin{column}{0.45\textwidth}
    \fcolorbox{Rule}{Paper}{%
  \begin{minipage}[t]{0.97\linewidth}
    \vspace{1.1mm}
    \hspace{1.4mm}\begin{minipage}{0.92\linewidth}
      \PanelTitle{内容说明}
      \scriptsize \begin{itemize}
  \item 用少量硬件扩展保护 attestation 代码、密钥和内存访问，并形式化验证整体 RA 安全属性
\end{itemize}
    \end{minipage}
    \vspace{1mm}
  \end{minipage}}
  \end{column}
  \begin{column}{0.49\textwidth}
    \fcolorbox{Rule}{Panel}{%
  \begin{minipage}[t]{0.97\linewidth}
    \vspace{1.1mm}
    \hspace{1.4mm}\begin{minipage}{0.92\linewidth}
      \PanelTitle{核心设计拆解}
      \scriptsize \begin{itemize}
  \item 01. Hardware/software RA root co-design
  \item 02. Formal verification of security sub-properties
  \item 03. MSP430/Basys3 FPGA implementation
\end{itemize}
    \end{minipage}
    \vspace{1mm}
  \end{minipage}}
  \end{column}
\end{columns}
\SourceFootnote{来源：VRASED verified RA co-design｜证据等级：E1 / 同行评审改进证据}
\end{frame}


\begin{frame}[t,shrink=6]{VRASED verified RA co-design：实验结果}
\DeckKicker{实验结果}
\SlideMeta{证据等级}{E1 / 同行评审改进证据}{来源状态：本地 PDF 已核验}
{\large\bfseries\color{Ink} USENIX Security 2019 peer-reviewed paper\par}\vspace{1.2mm}
\begin{columns}[T,totalwidth=\textwidth]
  \begin{column}{0.45\textwidth}
    \fcolorbox{Rule}{Paper}{%
  \begin{minipage}[t]{0.97\linewidth}
    \vspace{1.1mm}
    \hspace{1.4mm}\begin{minipage}{0.92\linewidth}
      \PanelTitle{内容说明}
      \scriptsize \begin{itemize}
  \item USENIX Security 2019 peer-reviewed paper
  \item reports prototype and verification evidence for simple embedded devices.
\end{itemize}
    \end{minipage}
    \vspace{1mm}
  \end{minipage}}
  \end{column}
  \begin{column}{0.49\textwidth}
    \fcolorbox{Rule}{Panel}{%
  \begin{minipage}[t]{0.97\linewidth}
    \vspace{1.1mm}
    \hspace{1.4mm}\begin{minipage}{0.92\linewidth}
      \PanelTitle{实验/证据边界}
      \scriptsize {\tiny
\begin{tabularx}{\linewidth}{@{}YYY@{}}
\toprule
\textbf{对象} & \textbf{可支撑} & \textbf{边界} \\
\midrule
数据/来源 & VRASED verified RA co-design & 本地 PDF 已核验 \\
Baseline/对照 & 以论文或规范原文为准 & 不跨材料外推 \\
指标/对象 & 只使用当前来源可证明的信息 & E1 / 同行评审改进证据 \\
使用边界 & 可进入本方向演进图 & 不能替代更强证据 \\
\bottomrule
\end{tabularx}
}
    \end{minipage}
    \vspace{1mm}
  \end{minipage}}
  \end{column}
\end{columns}
\SourceFootnote{来源：VRASED verified RA co-design｜证据等级：E1 / 同行评审改进证据}
\end{frame}


\begin{frame}[t,shrink=6]{VRASED verified RA co-design：文章评价}
\DeckKicker{文章评价}
\SlideMeta{证据等级}{E1 / 同行评审改进证据}{来源状态：本地 PDF 已核验}
{\large\bfseries\color{Ink} 证据强，结合形式化验证、硬件/软件 co-design 和真实原型\par}\vspace{1.2mm}
\begin{columns}[T,totalwidth=\textwidth]
  \begin{column}{0.45\textwidth}
    \fcolorbox{Rule}{Paper}{%
  \begin{minipage}[t]{0.97\linewidth}
    \vspace{1.1mm}
    \hspace{1.4mm}\begin{minipage}{0.92\linewidth}
      \PanelTitle{内容说明}
      \scriptsize \begin{itemize}
  \item 设计优势：证据强，结合形式化验证、硬件/软件 co-design 和真实原型。
  \item 局限性：目标是 simple embedded devices，不是 cloud CVM、CCA Realm 或 CoVE TVM evidence profile。
  \item 商业化潜力：适合可审计 MCU/IoT RoT 和 attestation IP；风险在 proof maintenance、toolchain 和 profile 标准化。
  \item 证据边界：E1 / Peer-reviewed SOTA；本地 PDF 已核验。
\end{itemize}
    \end{minipage}
    \vspace{1mm}
  \end{minipage}}
  \end{column}
  \begin{column}{0.49\textwidth}
    \fcolorbox{Rule}{Panel}{%
  \begin{minipage}[t]{0.97\linewidth}
    \vspace{1.1mm}
    \hspace{1.4mm}\begin{minipage}{0.92\linewidth}
      \PanelTitle{文章评价}
      \scriptsize {\tiny
\begin{tabularx}{\linewidth}{@{}YY@{}}
\toprule
\textbf{维度} & \textbf{判断} \\
\midrule
设计优势 & 证据强，结合形式化验证、硬件/软件 co-design 和真实原型。 \\
主要局限 & 目标是 simple embedded devices，不是 cloud CVM、CCA Realm 或 CoVE TVM evidence profile。 \\
商业落地 & 适合可审计 MCU/IoT RoT 和 attestation IP；风险在 proof maintenance、toolchain 和 profile 标准。 \\
讲稿定位 & 当前 SOTA/补充边界; 证据强度 E1 \\
\bottomrule
\end{tabularx}
}
    \end{minipage}
    \vspace{1mm}
  \end{minipage}}
  \end{column}
\end{columns}
\SourceFootnote{来源：VRASED verified RA co-design｜证据等级：E1 / 同行评审改进证据}
\end{frame}


\begin{frame}[t,shrink=6]{Attestation、boot、lifecycle：方向总结}
\DeckKicker{方向总结}
\SlideMeta{证据等级}{方向综述}{来源状态：公开材料综合}
{\large\bfseries\color{Ink} Attestation、boot、lifecycle 的结论是：三篇主讲材料共同给出技术演进，但仍要用证据等级约束每个结论。\par}\vspace{1.2mm}
\begin{columns}[T,totalwidth=\textwidth]
  \begin{column}{0.45\textwidth}
    \fcolorbox{Rule}{Paper}{%
  \begin{minipage}[t]{0.97\linewidth}
    \vspace{1.1mm}
    \hspace{1.4mm}\begin{minipage}{0.92\linewidth}
      \PanelTitle{内容说明}
      \scriptsize \begin{itemize}
  \item SWATT: 软件 attestation 早期代表：开创/基础入口，E1 / 基础证据。
  \item RATS 架构：代表性改进，E0 / Spec-standard SOTA。
  \item VRASED verified RA co-design：当前 SOTA/补充边界，E1 / 同行评审改进证据。
  \item 本方向结论：先用三篇主讲材料建立演进关系，再用证据等级控制机制、实验和商业化结论。
\end{itemize}
    \end{minipage}
    \vspace{1mm}
  \end{minipage}}
  \end{column}
  \begin{column}{0.49\textwidth}
    \fcolorbox{Rule}{Panel}{%
  \begin{minipage}[t]{0.97\linewidth}
    \vspace{1.1mm}
    \hspace{1.4mm}\begin{minipage}{0.92\linewidth}
      \PanelTitle{本方向方法对比}
      \scriptsize {\tiny
\begin{tabularx}{\linewidth}{@{}YYY@{}}
\toprule
\textbf{论文} & \textbf{定位} & \textbf{选择理由} \\
\midrule
SWATT: 软件 attestation 早期代表 & 开创/基础入口 & SWATT is the classic starting point for remote software attestation... \\
RATS 架构 & 代表性改进 & RFC 9334 defines the current RATS architecture for cross-TEE remote... \\
VRASED verified RA co-design & 当前 SOTA/补充边界 & VRASED is the strongest peer-reviewed verified RA co-design anchor for... \\
\bottomrule
\end{tabularx}
}
    \end{minipage}
    \vspace{1mm}
  \end{minipage}}
  \end{column}
\end{columns}
\SourceFootnote{来源：论文原文、官方规范或公开材料｜证据等级：见本页 badge}
\end{frame}
